% ==============================================================================
% Introduction Generale - Rapport PFE - ISET Nabeul
% ==============================================================================

\chapter*{Introduction g\'{e}n\'{e}rale}
\addcontentsline{toc}{chapter}{Introduction g\'{e}n\'{e}rale}
\markboth{Introduction g\'{e}n\'{e}rale}{Introduction g\'{e}n\'{e}rale}

\vspace{0.5cm}

% ---------------------------------------------------------------------------
% Contexte general : transformation numerique et livraison
% ---------------------------------------------------------------------------

La transformation num\'{e}rique constitue aujourd'hui un levier strat\'{e}gique
pour les entreprises de tous les secteurs. Dans le domaine de la gestion
commerciale et de la livraison, cette transformation se traduit par le
d\'{e}ploiement de solutions logicielles de plus en plus sophistiqu\'{e}es,
capables de g\'{e}rer l'ensemble de la cha\^{i}ne de valeur~: de la prise de
commande \`{a} la facturation, en passant par le suivi en temps r\'{e}el des
livraisons et la gestion des stocks.

Le march\'{e} tunisien, \`{a} l'instar des march\'{e}s internationaux, conna\^{i}t une
croissance significative de la demande en solutions de gestion int\'{e}gr\'{e}es.
Les entreprises cherchent \`{a} optimiser leurs processus op\'{e}rationnels, \`{a}
r\'{e}duire les erreurs humaines et \`{a} am\'{e}liorer la satisfaction de leurs
clients. Cette tendance s'accompagne d'exigences croissantes en mati\`{e}re de
fiabilit\'{e}, de performance et de qualit\'{e} des applications d\'{e}ploy\'{e}es.

\vspace{0.3cm}

% ---------------------------------------------------------------------------
% Presentation de DevWise et de ses projets
% ---------------------------------------------------------------------------

C'est dans ce contexte que s'inscrit l'activit\'{e} de \textbf{DevWise}, une
entreprise bas\'{e}e \`{a} Nabeul sp\'{e}cialis\'{e}e dans le d\'{e}veloppement de solutions
logicielles innovantes. DevWise con\c{c}oit et maintient deux produits
phares~:

\begin{itemize}
  \item \textbf{Speed Line}~: une application de gestion commerciale
    compl\`{e}te permettant aux entreprises de g\'{e}rer leurs ventes, leurs
    achats, leur stock, leur facturation et leurs tableaux de bord
    d\'{e}cisionnels. Speed Line se d\'{e}cline en version web et bureau, offrant
    une flexibilit\'{e} d'utilisation adapt\'{e}e aux besoins de chaque
    utilisateur.

  \item \textbf{Delivery App}~: un \'{e}cosyst\`{e}me de quatre applications
    (web et mobile) d\'{e}di\'{e}es \`{a} la gestion du cycle complet de livraison.
    Cet \'{e}cosyst\`{e}me couvre les besoins des diff\'{e}rents acteurs~: les
    administrateurs pour le pilotage global, les livreurs pour le suivi
    terrain, les clients pour le passage et le suivi de commandes, ainsi
    qu'une interface de gestion des restaurants et commerces partenaires.
\end{itemize}

Ces deux produits, en \'{e}volution constante, repr\'{e}sentent un p\'{e}rim\`{e}tre
fonctionnel riche et complexe. Chaque nouvelle version apporte son lot de
fonctionnalit\'{e}s, de corrections et d'am\'{e}liorations, ce qui multiplie les
sc\'{e}narios \`{a} valider avant chaque mise en production.

\vspace{0.3cm}

% ---------------------------------------------------------------------------
% Problematique : qualite logicielle et tests
% ---------------------------------------------------------------------------

Face \`{a} cette complexit\'{e} croissante, la qualit\'{e} logicielle devient un
enjeu majeur. Les tests manuels, bien qu'indispensables dans certaines
phases du d\'{e}veloppement, pr\'{e}sentent des limites significatives lorsqu'ils
constituent l'unique m\'{e}thode de validation~:

\begin{itemize}
  \item \textbf{Co\^{u}t en temps}~: l'ex\'{e}cution manuelle de l'ensemble des
    sc\'{e}narios de test devient de plus en plus chronophage \`{a} mesure que
    l'application s'\'{e}toffe. Un cycle de test complet peut n\'{e}cessiter
    plusieurs jours de travail.

  \item \textbf{Risque d'erreur humaine}~: la r\'{e}p\'{e}tition de t\^{a}ches
    identiques favorise l'inattention et les oublis. Certains cas limites
    ou r\'{e}gressions peuvent passer inaper\c{c}us.

  \item \textbf{Fr\'{e}quence de test limit\'{e}e}~: le co\^{u}t des tests manuels
    emp\^{e}che leur ex\'{e}cution syst\'{e}matique \`{a} chaque modification du code
    source, retardant ainsi la d\'{e}tection des anomalies.

  \item \textbf{Difficult\'{e} de mise \`{a} l'\'{e}chelle}~: avec la multiplication
    des plateformes (web, Android, iOS) et des configurations, la
    couverture exhaustive devient pratiquement impossible par des moyens
    exclusivement manuels.
\end{itemize}

L'automatisation des tests appara\^{i}t donc comme une r\'{e}ponse adapt\'{e}e \`{a}
ces d\'{e}fis. En mettant en place un robot de tests automatis\'{e}s, il devient
possible d'ex\'{e}cuter rapidement et de mani\`{e}re r\'{e}p\'{e}table un grand nombre de
sc\'{e}narios, d'identifier les r\'{e}gressions d\`{e}s leur apparition et de
lib\'{e}rer du temps pour les \'{e}quipes de d\'{e}veloppement, qui peuvent alors se
concentrer sur des t\^{a}ches \`{a} plus forte valeur ajout\'{e}e.

\vspace{0.3cm}

% ---------------------------------------------------------------------------
% Objectifs du stage
% ---------------------------------------------------------------------------

C'est dans cette perspective que s'inscrit notre projet de fin d'\'{e}tudes,
r\'{e}alis\'{e} au sein de l'entreprise DevWise sur une p\'{e}riode de quatre mois,
sous l'encadrement de M.~Mahmoud Slama. Les objectifs principaux de ce
stage sont les suivants~:

\begin{enumerate}
  \item \textbf{\'{E}tudier l'existant}~: analyser l'architecture des projets
    Speed Line et Delivery App, identifier les fonctionnalit\'{e}s critiques
    et recenser les sc\'{e}narios de test prioritaires.

  \item \textbf{Concevoir une architecture de tests}~: d\'{e}finir une
    architecture modulaire et maintenable pour le robot de tests, en
    s\'{e}lectionnant les outils et frameworks les plus adapt\'{e}s au contexte
    technique des projets.

  \item \textbf{Impl\'{e}menter les tests automatis\'{e}s}~: d\'{e}velopper les
    scripts de tests couvrant les sc\'{e}narios fonctionnels principaux des
    deux produits, en adoptant les bonnes pratiques du \textit{test
    automation}.

  \item \textbf{Int\'{e}grer les tests dans le processus de d\'{e}veloppement}~:
    mettre en place un pipeline d'int\'{e}gration continue permettant
    l'ex\'{e}cution automatique des tests \`{a} chaque modification du code.

  \item \textbf{Produire des rapports de tests}~: g\'{e}n\'{e}rer des rapports
    d\'{e}taill\'{e}s et exploitables pour faciliter le suivi de la qualit\'{e} et
    la prise de d\'{e}cision.
\end{enumerate}

\vspace{0.3cm}

% ---------------------------------------------------------------------------
% Methodologie adoptee
% ---------------------------------------------------------------------------

Pour mener \`{a} bien ce projet, nous avons adopt\'{e} une \textbf{m\'{e}thodologie
agile} inspir\'{e}e du framework \textbf{Scrum}. Cette approche it\'{e}rative et
incr\'{e}mentale nous a permis de~:

\begin{itemize}
  \item D\'{e}couper le projet en \textit{sprints} de deux semaines, chacun
    d\'{e}livrant un ensemble coh\'{e}rent de tests automatis\'{e}s.
  \item Organiser des r\'{e}unions r\'{e}guli\`{e}res avec l'encadrant pour
    synchroniser les priorit\'{e}s et ajuster le plan de travail.
  \item Int\'{e}grer les retours d'\'{e}quipe de mani\`{e}re continue pour
    am\'{e}liorer la qualit\'{e} des livrables.
  \item Maintenir un backlog prioris\'{e} des sc\'{e}narios \`{a} automatiser, en
    fonction de leur criticit\'{e} et de leur fr\'{e}quence d'utilisation.
\end{itemize}

Cette m\'{e}thodologie s'est av\'{e}r\'{e}e particuli\`{e}rement adapt\'{e}e au contexte du
projet, o\`{u} les exigences pouvaient \'{e}voluer en fonction des nouvelles
versions des applications et des retours des utilisateurs.

\vspace{0.3cm}

% ---------------------------------------------------------------------------
% Organisation du rapport
% ---------------------------------------------------------------------------

Le pr\'{e}sent rapport est structur\'{e} en plusieurs chapitres, organis\'{e}s de
mani\`{e}re \`{a} retracer la d\'{e}marche suivie tout au long de ce stage~:

\begin{itemize}
  \item \textbf{Chapitre 1~: Pr\'{e}sentation de l'organisme d'accueil et
    cadre du projet.} Ce chapitre pr\'{e}sente l'entreprise DevWise, son
    domaine d'activit\'{e}, les projets Speed Line et Delivery App, ainsi
    que le cadre g\'{e}n\'{e}ral et la probl\'{e}matique du stage.

  \item \textbf{Chapitre 2~: \'{E}tat de l'art et \'{e}tude th\'{e}orique.} Nous y
    passons en revue les concepts fondamentaux li\'{e}s aux tests logiciels,
    les diff\'{e}rentes strat\'{e}gies d'automatisation, ainsi qu'une \'{e}tude
    comparative des outils et frameworks disponibles sur le march\'{e}.

  \item \textbf{Chapitre 3~: Analyse et sp\'{e}cification des besoins.} Ce
    chapitre d\'{e}taille l'analyse des besoins fonctionnels et non
    fonctionnels, la d\'{e}finition des sc\'{e}narios de test prioritaires et
    la mod\'{e}lisation UML associ\'{e}e.

  \item \textbf{Chapitre 4~: Conception et architecture de la solution.}
    Nous y pr\'{e}sentons l'architecture technique retenue pour le robot de
    tests, les choix technologiques effectu\'{e}s et la conception d\'{e}taill\'{e}e
    des diff\'{e}rents modules.

  \item \textbf{Chapitre 5~: R\'{e}alisation et mise en \oe uvre.} Ce dernier
    chapitre couvre l'impl\'{e}mentation des tests automatis\'{e}s, leur
    int\'{e}gration dans le pipeline CI/CD, les r\'{e}sultats obtenus et les
    captures d'\'{e}cran illustrant le fonctionnement du robot.
\end{itemize}

\vspace{0.3cm}

Enfin, nous conclurons ce rapport par un bilan g\'{e}n\'{e}ral du stage, en
mettant en \'{e}vidence les comp\'{e}tences acquises, les difficult\'{e}s
rencontr\'{e}es et les perspectives d'\'{e}volution du projet.

\newpage
